\section{Experimental Design}
\label{sec:design}

\subsection{Disruption scenarios}

Closures are selected \emph{relative to the current optimal path}, since
removing random distant edges would change nothing. Table~\ref{tab:scenarios}
lists the seven conditions.

\begin{table}[t]
\caption{Disruption scenarios. Rejection rate is the share of candidate
closures discarded for disconnecting the OD pair.}
\label{tab:scenarios}
\centering
\begin{tabular}{llr}
\toprule
ID & Closure & Rejection \\
\midrule
S0 & None (routing-quality control) & --- \\
S1 & One random edge on the optimal path & 2.6\% \\
S2 & Highest-betweenness edge on the path & 1.0\% \\
S3 & Three simultaneous path edges & 14.5\% \\
S4 & Five simultaneous path edges & 31.7\% \\
S5 & Corridor: one named OSM way & 7.4\% \\
S6 & Intersection: a node and all incident edges & 4.5\% \\
\bottomrule
\end{tabular}
\end{table}

The closure unit is the ordered node pair, and all parallel edges on it close,
including the reverse direction. S5 closes one chosen way from the seed
edge's way-ID list rather than their union, since the union would close two
to four unrelated corridors and make severity depend on how OSMnx happened to
simplify the geometry. Source and target must remain connected after
removal; rejected candidates are resampled and \emph{counted}, because the
rejection rate is itself a measurement. Seeds derive from SHA-256 rather than
Python's built-in hash, which is randomized per process and would have made
the sweep irreproducible across parallel workers.

\subsection{Matrix and pairing}

Thirty origin--destination pairs are drawn with a fixed seed, stratified
10/10/10 by straight-line distance, and verified to have zero overlap with
the pairs used for hyperparameter tuning or for the ACO exit-criterion
measurement. The full matrix is 30 OD pairs $\times$ 7 scenarios $\times$ 2
traffic levels $\times$ \{warm, cold\}, giving 780 cells and
\textbf{5{,}520 runs} (ACO 3{,}900, Physarum 780, Dijkstra 420, A* 420).

ACO receives five seeds per cell. Physarum receives one, and this is
\emph{measured rather than assumed}: conductivities initialise from a fixed
value, and five OD pairs at three initialisation seeds produced costs
identical to four decimal places. Classical routers run cold only, since a
one-shot router has no warm mode.

Crucially, all three algorithms within one cell share \emph{one} Dijkstra
denominator, computed on the same weighted graph, so stretch is comparable
across algorithms within a cell by construction.

\subsection{Metrics}

Each run emits a flat record with its algorithm, OD pair, scenario, traffic
level, start mode, seed, config hash and timestamp. The reported quantities
are:

\begin{itemize}\setlength{\itemsep}{1pt}
\item \emph{Stretch}, $c_{\text{alg}} / c_{\text{Dijkstra}}$ on the same
  weighted graph. Values below 1.0 are a bug and are asserted against.
\item \emph{Post-disruption stretch}, against Dijkstra recomputed on the
  disrupted graph.
\item \emph{Recovery iterations} and \emph{recovery wall-clock}.
\item \emph{Walk-only stretch}, pre- and post-disruption: the maximum-flux
  walk's own result, taken against the same per-cell denominator as the
  pooled figures so the two are comparable within a row.
\item \emph{Extraction provenance}, the stage that produced the reported
  path, recorded at the iteration where the best cost improved.
\item \emph{Failure}, recorded and reported, never silently dropped.
\end{itemize}

One planned metric is excluded. \texttt{time\_to\_first\_valid\_path} was
intended to separate ``finds a valid route quickly'' from ``finds a good one
slowly'', but it is \textbf{exactly 1.0 on all 5{,}012 disrupted rows, for
every algorithm, with zero variance}. The cause is structural: ACO's ants
construct complete paths from the first iteration, and Physarum's extraction
falls back to Dijkstra, which always returns a path. It discriminates
nothing here and we exclude it rather than report a degenerate column. A
meaningful version would have to be redefined, for example as iterations to a first path within
$x\%$ of optimal, which is a different measurement.

\subsection{Statistics}

Distributions of recovery time and iterations are strongly right-skewed and,
for Physarum under warm start, bimodal. We therefore report \textbf{median
and interquartile range throughout, never means}, and compare with
\textbf{paired} tests: the Wilcoxon signed-rank
test~\cite{wilcoxon1945individual} on identical (OD pair, scenario, traffic
level, seed) tuples, with exact zero differences dropped as the statistic
requires. Confidence intervals come from a two-stage clustered bootstrap
(10{,}000 resamples): OD pairs are resampled first, then seeds within each
drawn pair, because rows from one OD pair are not independent.

Pre-disruption quality is deduplicated on the initial solve, which is
identical across all seven scenarios and both start modes; pooling raw rows
would count one solve up to thirteen times and shrink every interval by
bookkeeping alone.

Two rules follow from errors we made and corrected during this work, and we
state them because they generalise. \textbf{Never compare two medians whose
sample sizes differ without first checking that they cover the same cells}:
the design is fully paired by construction, so an unpaired comparison is
always a choice, and in our case it reversed the sign of a reported effect.
\textbf{Never report a bimodal distribution through a ratio of medians}: an
earlier version of our warm-start result did so and produced a spurious
``24$\times$'' speed-up.
